\documentclass[12pt]{article}
\usepackage[english]{babel}
\usepackage{graphicx}
\usepackage{framed}
\usepackage[normalem]{ulem}
\usepackage{amsmath}
\usepackage{amsthm}
\usepackage{amssymb}
\usepackage{amsfonts}
\usepackage{enumerate}
\usepackage{graphicx} % Allows including image
\usepackage{hyperref}
\hypersetup{
  linkcolor  = blue,
  citecolor  = blue,
  urlcolor   = blue,
  colorlinks = true,
} % color setup
\usepackage[utf8]{inputenc}
\usepackage[margin=1in]{geometry}
\usepackage{float} % for "H" placement specifier
\usepackage[authoryear]{natbib}
\usepackage{indentfirst}
\usepackage[dvipsnames]{xcolor}
\usepackage{float}
\usepackage{array}


\usepackage[T1]{fontenc}
\usepackage{tikz}
\newcommand{\fatterdot}{\raisebox{0.15ex}{\tikz\filldraw[black,x=1.5pt,y=1.5pt] (0,0) circle (1);}}
\newcommand{\xn}{x^{\frac{1}{n}}}
\newcommand{\xnone}{x^{\frac{1}{n+1}}}

\newcommand{\R}{\mathbb{R}}
\DeclareUnicodeCharacter{FB01}{fi}
\DeclareUnicodeCharacter{FF0C}{,}
\DeclareMathOperator*{\Prob}{P}
\DeclareMathOperator*{\argmax}{arg\,max}
\DeclareMathOperator*{\argmin}{arg\,min}
\setlength\extrarowheight{1pt}

\newcommand\red[1]{{\color{red}#1}}

\bibliographystyle{plainnat}

\newtheorem{theorem}{Theorem}[section]
\newtheorem{corollary}{Corollary}[theorem]
\newtheorem{lemma}[theorem]{Lemma}
\newtheorem{prop}[theorem]{Proposition}
\newtheorem{defn}[theorem]{Definition}
\allowdisplaybreaks

\title{An Investigation into Paradoxes Between Latent Aptitude and Observed Statistics via Novel Full House Modeling}
\author{Daniel J. Eck$^1$, Shen Yan$^1$}

\date{}

\begin{document}

\maketitle

{\centering
1. Department of Statistics, University of Illinois Urbana-Champaign \\}


\vspace{12pt}

\noindent This is a hopefully self-contained writeup that contains the foundational material needed to understand what results we need to move this work forward. 

\vspace{12pt}

\noindent \textbf{Sentences in bold text will denote results that we need investigated.}

\vspace{12pt}

\noindent \textcolor{red}{Sentences in red text will denote attempts that we have made.}

\section*{Summary}

\begin{itemize}
    \item \textbf{Contribution of the paper:}
    \begin{itemize}
        \item Introduces a statistical framework (Full House Modeling) to fairly compare individuals across contexts (schools, eras, and other domains) where measurement and selection differ. 
        \item This work formalizes the model that we developed in our paper on comparing baseball players across eras \citep{yan2025comparing}.
        \item Defines and quantifies the \emph{aptitude--achievement paradox}: cases where someone with higher true ability shows lower observed performance across contexts. We define this paradox in terms of representative individuals from two systems drawn from two distinct eligible populations $i$. For example, this could be government officials elected from two states or baseball players from different seasons.
        \item Derives \emph{distribution-free formulas} for paradox probabilities, connects to classic paradoxes (Simpson's, Lord's), and provides modern fairness tools.
    \end{itemize}

    \item \textbf{What we are trying to solve now (generalization of Proposition 3.5):}
    \begin{itemize}
        \item Extend the proof of the key \emph{inequality condition} that underpins the distribution-free paradox probability formula.
        \item Notation: $N_i$ is the total size of population $i$, and $n_i$ is the number of top individuals from that population selected into the system (so $n_i \leq N_i$).
        \item The condition is currently proven when $n_1 \leq n_2$ and $N_2 \leq N_1$.
        \item Goal is to generalize it to the following condition:
        \[
            \frac{n_1}{N_1} \;\leq\; \frac{n_2}{N_2},
        \]
        i.e., system 1 selects a smaller fraction of its population than system 2. We have repeatedly verified this result empirically, but cannot yet prove it        
        \item Solving this unlocks the full generality of our paradox probability results for top-selection settings to the best of our knowledge.
    \end{itemize}
\end{itemize}



What follows is the current version of the paper up to Proposition 3.5. Proofs of results are included for completeness.

\section*{Introduction}

The following paragraph introduces the intuition behind Full House Modeling. It is adapted from \cite{yan2025comparing} with some slight changes: 

\vspace{12pt}

Comparisons of human performance across differing contexts, schools, regions, eras, institutions, and tasks, are central to education, labor economics, psychometrics, medicine, and sports analytics \citep{chetty2014measuring, iezzoni1997risk, bartholomew2011latent, bond2013applying, BRL99, imbens2018causal, holland2012differential, yan2025comparing}. A recurring challenge is that an individual in one context may have higher latent aptitude yet lower observed achievement than a counterpart elsewhere. We term this the aptitude-achievement paradox and develop a statistical framework to quantify its probability across contexts. Formally, let individual $j$ in context $i$ have an observed $\R$-valued statistic $Y_{i,j}$ and an underlying latent $\R$-valued aptitude $X_{i,j}$, and let individual $n$ in context $m$ has an observed $\R$-valued statistic $Y_{m,n}$,and an underlying latent $\R$-valued aptitude $X_{m,n}$. The probability of aptitude-achievement paradox is 
$$
\Prob\left(X_{i,j} > X_{m,n}, Y_{i,j} < Y_{m,n}\right)
$$

The paradox is structurally impossible if a single, strictly increasing, context-invariant measurement map aptitude to achievement and if selection into tasks is ignorable. In practice, however, both measurement heterogeneity (non-invariance, differential item functioning) and non-ignorable selection (who attempts which tasks, faces which opponents, receives which treatment) are pervasive. These features can reverse cross-context rankings even when the latent ranking is unambiguous, echoing classic reversal phenomena such as Simpson's \citep{simpson1951interpretation} and Lord's paradoxes \citep{lord1967paradox} while implicating modern concerns about fairness, transportability, and domain shift. What practitioners lack is a principled way to quantify how often such reversals should occur under an explicit model of aptitude, measurement, and selection.

A large class of methods measures how often observed orderings disagree with a target ordering-precisely the phenomenon at the heart of our cross-context paradox. At the most elementary level, discordance statistics summarize the share of ordered pairs whose observed outcomes align (or misalign) with a reference score or latent ordering. Classic measures include Kendall's $\tau$ \citep{kendall1938new} and Somers'D \citep{somers1962new}, which compute differences between the proportions of concordant and discordant pairs and offer adjustments for ties; Somers'D is asymmetric and treats one variable (e.g., a score) as the predictor of the other (e.g., an outcome), making it particularly natural for "does the score rank the outcomes correctly?" questions. These rank-based indices provide assumption-lean, effect-size-like summaries of reversal risk. However, because these indices operate on observed rankings, they neither place individuals on a common latent scale nor disentangle measurement non-invariance or selection, motivating models that make the measurement process explicit.

Traditional latent trait models \citep{bartholomew2011latent} focus on estimating underlying traits from observed categorical data. One of the example of traditional latent trait models, the Rasch model \citep{bond2013applying} ensures that comparisons of individuals' abilities are independent of the specific items administered, and the Rasch model expect no differential item functioning (DIF) \citep{holland2012differential}, which means an item behaves differently for two (or more) groups of examinees who have the same underlying ability on the target construct. But it does not explicitly leverage the source population structure. Likewise, the bayesian hierarchical modeling \citep{BRL99} with the application in comparison of individuals' latent aptitudes typically captures population-level variation through seasonal random effects. In parallel, the causal bootstrap \citep{imbens2018causal} highlights that valid causal inference similarly requires situating the observed sample within large source population to properly account for sampling uncertainty. 

Full House Modeling offers a novel approach for comparing individuals across diverse contexts, initially developed through the analysis of baseball players across different eras \citep{yan2025comparing}. This analytical framework gained further historical significance \citep{burgos2025change}, which was subsequently explored and presented in a separate study. These analyses are transforming baseball analytics and are increasingly recognized within the broader baseball community. The applicability of Full House Modeling was further extended to investigate the influence of historical figures \citep{cohen2025contextualizing}. 

In this article, we formalize the Full House Methodology and investigate the probabilities of observing paradoxes between latent aptitude and observed achievement in comparative analysis of individuals across contexts. The aptitude-achievement  paradox arises when an individual from one context demonstrates higher latent aptitude but simultaneously lower observed achievement than an individual from another context. We provide explicit probabilities for these paradoxes in two scenarios: first, when observed achievement is independent and identically distributed (i.i.d.), and second, when observed statistics across contexts follow different normal distributions but remain independent. We also explore bounds for these paradoxes in more generalized settings and demonstrate that when observed achievement is normally distributed, paradox probabilities can be reliably estimated using a plug-in method.


\section*{Model Setting}

We now motivate the structure of the Full House Model. Let $N_i$ denote the size of population $i$. We will suppose the every individual $k = 1,\ldots, N_i$ has an underlying latent $\R$-valued aptitude 
$X_{i,1}, \ldots, X_{i,N_i} \overset{iid}{\sim} F_{X}$. Let $\textbf{X}_{i,-j}$ be the vector of all latent aptitudes not including the component $j$. We then define an indicator function $1_{i}\left(X_{i,j},\textbf{X}_{i,-j}\right)$ as a determined and known system inclusion function for population $i$ where a 1 indicates inclusion and a 0 indicates exclusion. 
%$g_{i}\left(X_{i,j}, \textbf{X}_{i,-j}\right)=1$ indicates that individual $j$ is included in the $i$th system, and $g_{i}\left(X_{i,j}, \textbf{X}_{i,-j}\right)=0$ indicates that individual $j$ is not included in the $i$th system. 
We now suppose that an individual $j$ included in system $i$ has an observed $\R$-valued statistic $Y_{i,j} \overset{iid}{\sim} F_{Y_i}$, $j = 1,\ldots,n_i$, $n_i \leq N_i$. As we will see later, the observed $Y_{i,j}$ values are one-to-one realizations of $X_{i,j}$, provided that $g_{i}\left(X_{i,j}, \textbf{X}_{i,-j}\right)=1$. Therefore $Y_{i,1}, \ldots, Y_{i,n_i} \overset{iid}{\sim} F_{Y_i}$. 

Let $X_{i,(j)}$ and $Y_{i,(j)}$ denote the $j$th order statistic for system $i$. Here, larger numbers will correspond to a notion of "better." Thus $X_{i,(j)} > X_{i,(k)}$ whenever $j>k$, for example. We now form $(Y,X)$ pairs of random variables for each individual included in the system. These pairs are of the form $(Y_{i,(j)}, X_{i,(r_{i,j})})$ where $r_{i,j}$ indexes the aptitude of the $j$th highest value observed from the system. As an example, let $1_{i}\left(X_{i,j}, \textbf{X}_{i,-j}\right) = 1(X_{i,j} \geq X_{i,(N_i-n_i+1)})$, then the $(Y_{i,(j)}, X_{i,(r_{i,j})})$ pairs take the form $(Y_{i,(j)}, X_{i,(N_i - n_i + j)})$ for $j = 1,\ldots,n_i$.

%where we suppose that $Y_{i,1}, \ldots, Y_{i,n_i} \overset{iid}{\sim} F_{Y_i}$, $F_{Y_i}$ is continuous, and $n_i < N_i$ is the number of individuals included in system $i$. We denote $Y_{i,(j)}$ as the $j$th ordered statistic for system $i$.
The key Full House Modeling is the explicit connection between the elements forming the $(Y_{i,(j)}, X_{i,(r_{i,j})})$ pairs. Full House Modeling states this connection as
\begin{equation} \label{eq:truepar}
\begin{split}
  Y_{i, (j)} &= F_{Y_{i, (j)}}^{-1} \circ F_{X_{i, (r_{i,j})}} \left(X_{i, (r_{i,j})}\right), \\
  X_{i, (r_{i,j})} &= F_{X_{i, (r_{i,j})}}^{-1} \circ F_{Y_{i,(j)}} \left(Y_{i,(j)} \right),
\end{split}
\end{equation}
where $1_i(X_{i, (r_{i,j})}, \textbf{X}_{i, -(r_{i,j})}) = 1$, and both $F_{Y_i}$ and $F_X$ are continuous and invertible. We summarise Full House Modeling in the definition below.


\begin{defn} [Full House Model]
\label{def:fhm}
	Consider populations with sizes $N_i$ and system sizes $n_i$, $N_i \geq n_i$. Let a system inclusion function $1_i(\cdot,\cdot)$ be specified. Let $X_{i,j} \overset{iid}{\sim} F_X$, $j = 1,\ldots,N_i$. Let $Y_{i,j} \overset{iid}{\sim} F_{Y_i}$, $j = 1,\ldots, n_i$. Form pairs $(Y_{i,(j)}, X_{i,(r_{i,j})})$, $j = 1,\ldots,n_i$ for all individuals such that $1_i(X_{i, (r_{i,j})}, \textbf{X}_{i, -(r_{i,j})}) = 1$. Then the Full House Model generates $Y_{i,(j)}$ values following \eqref{eq:truepar} where $F_{Y_i}$ and $F_X$ are continuous and invertible.
\end{defn}

Note that Definition~\ref{def:fhm} fixes $F_X$ but allows $F_{Y_i}$ to vary across systems. 

%\vspace{12pt}\textbf{Remark}: The definition of a Full House Model immediately implies structural missingness. Every $X_{i,j}$ such that $1_i(X_{i,j},\textbf{X}_{i,-j}) = 0$ will not have a corresponding $Y$ value. In the following sections we will consider probabilities involving both $X$ and $Y$ quantities. To handle this structural missingness, we will suppose that $Y_{i,j}$ values that are not observed are not $\Prob$-measurable.
%\item The probability integral transformation steps in \eqref{eq:truepar} involving $F_{X_{i, (r_{i,j})}} \left(X_{i, (r_{i,j})}\right)$ and $F_{Y_{i,(j)}} \left(Y_{i,(j)} \right)$ may be cumbersome to evaluate in practice. Instead one may replace $F_{X_{i, (r_{i,j})}} \left(X_{i, (r_{i,j})}\right)$ with $F_{U_{i,(r_{i,j})}} \circ F_X \left(X_{i, (r_{i,j})}\right)$ where $U_{i,(j)} \sim $ \,Beta$(r_{i,j},N_i-r_{i,j}+1)$, and may replace $F_{Y_{i,(j)}} \left(Y_{i,(j)} \right)$ with $F_{U_{i,(j)}} \circ F_{Y_i}\left(Y_{i,(j)} \right)$ where $U_{i,(j)} \sim $\, Beta$(j,n_i-j+1)$.







\section*{Paradoxes Under Identical Distributions}
\label{sec:samedist}

In this section we will provide explicit paradox probability results in a Full House Modeling setting involving two populations with the additional specification that $F_{Y_i} = F_{Y}$, $i=1,2$ in Definition~\ref{def:fhm}. The following definition formally defines the notion of paradox used throughout this article.

\begin{defn}\label{def:paradox}
Consider the Full House Model in Definition~\ref{def:fhm} with two populations denoted by $i=1,2$. We say a paradox has occurred when both $X_{1,(r_{1,j})} > X_{2,(r_{2,k})}$ and  $Y_{1,(j)} < Y_{2,(k)}$ have occurred. The paradox probability is defined as
\begin{equation} \label{eq:simpleparadox}
  \Prob\left(X_{1,(r_{1,j})} > X_{2,(r_{2,k})}, Y_{1,(j)} < Y_{2,(k)}\right).
\end{equation}	
\end{defn}

As a simple illustrative example of the above, let $r_{i,t} = N_i - n_i + t$ and set $j = n_1$ and $k = n_2$, and further suppose that $Y_{1,j} < Y_{2,k}$ a.s. Then the paradox probability \eqref{eq:simpleparadox} is equivalent to $\Prob\left(X_{1,(N_1)} > X_{2,(N_2)}\right)$, and we have $\Prob\left(X_{1,(N_1)} > X_{2,(N_2)}\right) \to 1$ when $N_1 \to \infty$. We now introduce a convenient representation which allows the paradox probability to be expressed in terms of the latent aptitudes alone.


\begin{lemma} \label{lemma:rep}
Consider the Full House Model in Definition~\ref{def:fhm} with two populations, where $F_{Y_i} = F_{Y}$ for $i=1,2$. Then, 
\begin{equation} \label{eq:pararep}
  \Prob\left(X_{1,(r_{1,j})} > X_{2,(r_{2,k})}, Y_{1,(j)} < Y_{2,(k)}\right) 
  = \Prob\left( X_{2,(r_{2,k})} < X_{1,(r_{1,j})}  < h_{1,2}\left(X_{2,(r_{2,k})}\right)\right),
\end{equation}
where $h_{1,2,j,k}(\cdot) = F_{X_{1,(r_{1,j})}}^{-1} \circ F_{Y_{1, (j)}} \circ F_{Y_{2, (k)}}^{-1} \circ F_{X_{2,(r_{2,k})}} (\cdot)$.
\end{lemma}

\begin{proof}
Define $f_{1,j}(\cdot) = F_{Y_{1, (j)}}^{-1} \circ F_{X_{1,(r_{1,j})}} (\cdot)$ and 
$f_{2,k}(\cdot) = F_{Y_{2, (k)}}^{-1} \circ F_{X_{2,(r_{2,k})}} (\cdot)$. Then we have 
\begin{align*}
  &\Prob\left(X_{1,(r_{1,j})} > X_{2,(r_{2,k})}, Y_{1,(j)} < Y_{2,(k)}\right) \\
  &\qquad= \Prob\left(X_{1,(r_{1,j})} > X_{2,(r_{2,k}))}, f_{1,j}(X_{1,(r_{1,j})}) < f_{2,k}(X_{2,(r_{2,k})})\right) \\
  &\qquad= \Prob\left( X_{2,(r_{2,k})} < X_{1,(r_{1,j})}  < h_{1,2}\left(X_{2,(r_{2,k})}\right)\right),
\end{align*}
where $h_{1,2,j,k}(\cdot) = (f_{1,j}^{-1} \circ f_{2,k})(\cdot)$. %Existence of $f_i$ follows directly from the existence of $F_Y^{-1}(y)$, and existence of $h_{i,j}$ follows directly from the existence of $f_i$ and the existence of $F_X^{-1}(x)$.
\end{proof}

Our first formal result provides conditions that allow for an explicit paradox probability result. Those conditions require the following definition.

\begin{defn}
\label{def:ribf}
Let $\alpha,\beta > 0$. Then the regularized incomplete beta function is defined as
$$
  I_{\alpha,\beta}(x) = \frac{1}{B(\alpha,\beta)} \int_0^x u^{\alpha-1}(1-u)^{\beta-1}du,
$$
where $B(\alpha,\beta) = \frac{\Gamma(\alpha)\Gamma(\beta)}{\Gamma(\alpha + \beta)}$ is the beta coefficient.
\end{defn}

Note that the regularized incomplete beta function is the cumulative distribution function of a $\text{Beta}(\alpha,\beta)$ distribution. 

\begin{theorem} \label{theorem:samedistgeneral}
Consider the Full House Model in Definition~\ref{def:fhm} with two populations, where $F_{Y_i} = F_{Y}$ for $i=1,2$. Denote by $I$ the regularized incomplete beta function. Further assume that
\begin{equation} \label{eq:technical-condition}
  I_{j, n_1-j+1} \circ I_{k, n_2-k+1}^{-1} (t) - I_{r_{1,j}, N_1 - r_{1,j}+1} \circ I_{r_{2,k}, N_2 - r_{2,k} + 1}^{-1}(t) \geq 0, \qquad \text{for all} \; t \in (0,1).
\end{equation}
Under these assumptions,
\begin{align*}
	&\Prob\left(X_{1,(r_{1,j})} > X_{2,(r_{2,k})}, Y_{1,(j)} < Y_{2,(k)}\right) \\
  &= \sum_{p=j}^{n_1}\left(\begin{array}{l}
n_1 \\
p
\end{array}\right)\frac{B(k+p, n_2-n_1-p+k+1)}{B(k, n_2-k+1)} - \sum_{p=r_{1,j}}^{N_1}\left(\begin{array}{l}
N_1 \\
p
\end{array}\right)\frac{B(r_{2,k}+p, N_2-N_1-p+r_{2,k}+1)}{B(r_{2,k}, N_2-r_{2,k}+1)}.
\end{align*}
\end{theorem}


\noindent \textbf{Remarks}:
\begin{enumerate}
	\item The result of Theorem~\ref{theorem:samedistgeneral} is distribution free. It only depends on the sample sizes $n_i, N_i$, $i = 1,2$, the system inclusion function, and the indexes of the order statistics under comparison. The distribution on $Y$ is eliminated by Lemma~\ref{lemma:rep}. To see this, consider the distribution $F_{Y_{(j)}}$ for some $j = 1,\ldots,n_1$. The random variable $Y_{(j)} = F_Y^{-1} \circ U_{(j)}$ where $U_{(j)} \sim \text{Beta}(j,n_1-j+1)$. Thus, $F_{Y_{(j)}}(y) = I_{j,n_1-j+1} \circ F_Y (y)$. With this in mind, from Lemma~\ref{lemma:rep} we have that
\begin{align*}
	h_{1,2,j,k}(x) &= F_{X_{1,(r_{1,j})}}^{-1} \circ F_{Y_{1, (j)}} \circ F_{Y_{2, (k)}}^{-1} \circ F_{X_{2,(r_{2,k})}} (x) \\
	&= F_{X_{1,(r_{1,j})}}^{-1} \circ I_{j,n_1-j+1} \circ \left(F_Y \circ F_Y^{-1}\right) \circ I^{-1}_{k,n_2-k+1} \circ F_{X_{2,(r_{2,k})}} (x) \\
	&= F_{X_{1,(r_{1,j})}}^{-1} \circ I_{j,n_1-j+1}  \circ I^{-1}_{k,n_2-k+1} \circ F_{X_{2,(r_{2,k})}} (x).
\end{align*}
The above is the upper-bound of an integral with respect to $X$s in the proof of Theorem~\ref{theorem:samedistgeneral}. The influence of the distributional specification on $Y$ does not influence the calculation of the paradox probability result in Theorem~\ref{theorem:samedistgeneral}.
	\item The distribution on $X$ is eliminated when calculating the paradox probability \eqref{eq:simpleparadox} in its Lemma~\ref{lemma:rep} representation, $\Prob\left( X_{2,(r_{2,k})} < X_{1,(r_{1,j})}  < h_{1,2}\left(X_{2,(r_{2,k})}\right)\right)$. These calculations involve double integrals with respect to densities for $F_{X_{1,(r_{1,j})}}$ and $F_{X_{2,(r_{2,k})}}$. One integral evaluation and a reparameterization remove the specification of $F_X$ from the probability calculation as seen in the proof of Theorem~\ref{theorem:samedistgeneral}.
	\item The steps mentioned in the last sentence of the previous remark allow the probability calculation for $\Prob\left( X_{2,(r_{2,k})} < X_{1,(r_{1,j})}  < h_{1,2}\left(X_{2,(r_{2,k})}\right)\right)$ to simplify to 
$$
	\int_0^1\left\{I_{j, n_1-j+1} \circ I_{k, n_2-k+1}^{-1} (t) - I_{r_{1,j}, N_1 - r_{1,j}+1} \circ I_{r_{2,k}, N_2 - r_{2,k} + 1}^{-1}(t)\right\}^+ dt.
$$
The technical condition \eqref{eq:technical-condition} ensures that this integral defines a non-negative probability in strong terms. We demonstrate with analytic and numeric arguments when \eqref{eq:technical-condition} is satisfied and the results of Theorem~\ref{theorem:samedistgeneral} therefore hold.
\end{enumerate}


\begin{proof}
First, Lemma~\ref{lemma:rep} allows us to write
$$
	\Prob\left(X_{1,(r_{1,j})} > X_{2,(r_{2,k})}, Y_{1,(j)} < Y_{2,(k)}\right) = \Prob\left( X_{2,(r_{2,k})} < X_{1,(r_{1,j})}  < h_{1,2,j,k}\left(X_{2,(r_{2,k})}\right)\right),
$$
where $h_{1,2,j,k}(\cdot) = F_{X_{1,(r_{1,j})}}^{-1} \circ F_{Y_{1, (j)}} \circ F_{Y_{2, (k)}}^{-1} \circ F_{X_{2,(r_{2,k})}} (\cdot)$. We can rewrite $h_{1,2,j,k}(\cdot)$ into a form that will allow for specific computation involving Beta distributions. o see this, consider the distribution $F_{Y_{(j)}}$ for some $j = 1,\ldots,n_1$. The random variable $Y_{(j)} = F_Y^{-1} \circ U_{(j)}$ where $U_{(j)} \sim \text{Beta}(j,n_1-j+1)$. Thus, $F_{Y_{(j)}}(y) = g_{n,j} \circ F_Y (y)$,
%\begin{align*}
%  P\left(X_{2,(r_{2,k})} < x\right) &= P\left(U_{(r_{2,k})} < F_X(x)\right) \\ 
%  &= \frac{1}{B(r_{2,k},N_{2}-r_{2,k}+1)}\int_0^{F_X(x)}u^{r_{2,k}-1}(1-u)^{N_2-r_{2,k}} du	\\
%  &= g_{N_2,r_{2,k}} \circ F_X (x),
%\end{align*}
where we write $g_{n,j}(x) = I_{j, n-j+1}(x)$ for ease of notation. This expression of distribution functions allows us to write 
\begin{align*}
	&\Prob\left( X_{2,(r_{2,k})} < X_{1,(r_{1,j})}  < h_{1,2,j,k}\left(X_{2,(r_{2,k})}\right)\right) \\
	&= \int_{-\infty}^{+\infty}\left\{\int_{x}^{F_X^{-1} \circ g_{N_1,r_{1,j}}^{-1} \circ g_{n_1, j} \circ g_{n_2,k}^{-1} \circ g_{N_2,r_{2,k}} \circ F_X(x)} f_{X_{1,(r_{1,j})}}(y)\ dy\right\} \ f_{X_{2,(r_{2,k})}}(x) \ dx \\
	&= \int_{-\infty}^{+\infty}\left\{\int_{x}^{F_X^{-1} \circ g_{N_1,r_{1,j}}^{-1} \circ g_{n_1, j} \circ g_{n_2,k}^{-1} \circ g_{N_2,r_{2,k}} \circ F_X(x)} d \ g_{N_1, r_{1,j}}(F_X(y))\right\} \ f_{X_{2,(r_{2,k})}}(x) \ dx \\
	&= \int_{-\infty}^{+\infty}\left\{g_{n_1, j} \circ g_{n_2, k}^{-1} \circ g_{N_2, r_{2,k}} \circ F_X(x) - g_{N_1, r_{1,j}} \circ F_X(x)\right\}^+ \ f_{X_{2,(r_{2,k})}}(x) \ dx \\
	&= \int_{-\infty}^{+\infty}\left\{g_{n_1, j} \circ g_{n_2, k}^{-1} \circ g_{N_2, r_{2,k}} \circ F_X(x) - g_{N_1, r_{1,j}} \circ F_X(x)\right\}^+ \  d \ g_{N_2, r_{2,k}}(F_X(x)) \\
	&= \int_0^1 \left\{g_{n_1, j} \circ g_{n_2, k}^{-1}(t) - g_{N_1, r_{1,j}} \circ g_{N_2, r_{2,k}}^{-1}(t)\right\}^+	 \ dt.
\end{align*}
Condition~\eqref{eq:technical-condition} states that 
$$
\int_0^1 \left\{g_{n_1, j} \circ g_{n_2, k}^{-1}(t) - g_{N_1, r_{1,j}} \circ g_{N_2, r_{2,k}}^{-1}(t)\right\}^+\ dt = 
\int_0^1 \left\{g_{n_1, j} \circ g_{n_2, k}^{-1}(t) - g_{N_1, r_{1,j}} \circ g_{N_2, r_{2,k}}^{-1}(t)\right\}\ dt.
$$
Thus,
\begin{align*}	
&\int_0^1 \left\{g_{n_1, j} \circ g_{n_2, k}^{-1}(t) - g_{N_1, r_{1,j}} \circ g_{N_2, r_{2,k}}^{-1}(t)\right\}^+\ dt \\
	&\qquad= \int_0^1 g_{n_1, j}(x) \ dg_{n_2, k}(x)- \int_0^1g_{N_1, r_{1,j}}(x) \ dg_{N_2, r_{2,k}}(x)\\
	&\qquad= \int_0^1 I_x(j, n_1 -j+1) \frac{x^{k-1}(1-x)^{n_2-k}}{B(k, n_2-k+1)} \ dx \\ 
	&\qquad\qquad- \int_0^1 I_x(r_{1,j}, N_1-r_{1,j}+1) \frac{x^{r_{2,k}-1}(1-x)^{N_2-r_{2,k}}}{B(k, N_2-k+1)} \ dx\\
	&\qquad =\int_0^1 \sum_{p=j}^{n_1}\left(\begin{array}{l}
n_1 \\
p
\end{array}\right) x^p(1-x)^{n_1-p} \frac{x^{k-1}(1-x)^{n_2-k}}{B(k, n_2-k+1)} \ dx \\
&\qquad\qquad - \int_0^1 \sum_{p=r_{1,j}}^{N_1}\left(\begin{array}{l}
N_1 \\
p
\end{array}\right) x^p(1-x)^{N_1-p} \frac{x^{r_{2,k}-1}(1-x)^{N_2-r_{2,k}}}{B(r_{2,k}, N_2-r_{2,k}+1)} \ dx\\ 
&\qquad = \sum_{p=j}^{n_1}\left(\begin{array}{l}
n_1 \\
p
\end{array}\right)\frac{B(k+p, n_2+n_1-p-k+1)}{B(j, n_2-k+1)} \\ 
&\qquad\qquad - \sum_{p=r_{1,j}}^{N_1}\left(\begin{array}{l}
N_1 \\
p
\end{array}\right)\frac{B(r_{2,k}+p, N_2+N_1-p-r_{2,k}+1)}{B(r_{2,k}, N_2-r_{2,k}+1)}.
\end{align*}
\end{proof}


Often it is of interest to compare the individuals with the highest aptitudes in their respective populations. Below we state the paradox probability in such setting using a Full House Model with two populations with the specification that $F_{Y_1} = F_Y$ for $i=1,2$.

\begin{corollary} \label{corollary:samedisttop}
Suppose that the conditions of Theorem~\ref{theorem:samedistgeneral} hold. Suppose that $1_1(X_{1,(N_1)}, \textbf{X}_{1,-(N_1)}) = 1_2(X_{2,(N_2)}, \textbf{X}_{2,-(N_2)}) = 1$. Then,
\begin{align*}
	&\Prob\left(X_{1,(N_1)} > X_{2,(N_2)}, Y_{1,(n_1)} < Y_{2,(n_2)}\right) 
  = \begin{cases}
      \frac{n_2}{n_1 + n_2} - \frac{N_2}{N_1 + N_2}, & \text{if}\ \frac{n_1}{N_1} \leq \frac{n_2}{N_2} \\
      0, & \text{otherwise}
    \end{cases}.
\end{align*}
\end{corollary}

A similar result holds for the individuals with the lowest latent aptitudes in the population.

\begin{corollary} \label{corollary:samedistbottom}
Suppose that the conditions of Theorem~\ref{theorem:samedistgeneral} hold. Suppose that $1_1(X_{1,(1)}, \textbf{X}_{1,-(1)}) = 1_2(X_{2,(1)}, \textbf{X}_{2,-(1)}) = 1$. Then,
\begin{align*}
	&\Prob\left(X_{1,(1)} > X_{2,(1)}, Y_{1,(1)} < Y_{2,(1)}\right) 
  = \begin{cases}
      \frac{N_2}{N_1 + N_2} - \frac{n_2}{n_1 + n_2}, & \text{if}\ \frac{n_1}{N_1} \geq \frac{n_2}{N_2} \\
      0, & \text{otherwise}
    \end{cases}.
\end{align*}
\end{corollary}


We now demonstrate a setting in which technical condition \eqref{eq:technical-condition} holds. This setting is based on a configuration of $n_j$, $N_j$, $j=1,2$ under the assumption that the systems are populated with individuals from the source populations who have the highest aptitudes.


\begin{prop} \label{prop:condition}
Consider the Full House Model in Definition~\ref{def:fhm} with two populations, where $F_{Y_i} = F_{Y}$ for $i=1,2$, and let the system inclusion function be $1_i(X_{i,j} \geq X_{i,(N_i - n_i + 1)})$. Denote by $I$ the regularized incomplete beta function. Suppose that $\frac{n_1}{n_2}\leq 1$ and $\frac{N_2}{N_1} \leq 1$, and set $k = j$. Under these assumptions, we conclude that technical condition \eqref{eq:technical-condition} holds and the conclusion of Theorem~\ref{theorem:samedistgeneral} applies.	
\end{prop}

\begin{proof}


Specification of $1_i(X_{i,j} \geq X_{i,(N_i - n_i + 1)})$ implies that $r_{i,j} = N_i + 1 - j$. 
\vspace{12pt}

When either $n_1 = N_1$ and $n_2 = N_2$ or $n_1 = n_2$ and $N_1 = N_2$ we have, for all $x \in (0,1)$, that
$$
  I_{N_2-j,j+1} \circ I_{N_1-j,j+1}^{-1} \circ I_{n_1-j,j+1} \circ I_{n_2-j,j+1}^{-1} (t) \geq t.
$$
The left hand side of the above equation increases whenever $n_1$ or $N_1$ are increased \citep{askitis2021logarithmic}, or $n_2$ or $N_1$ are decreased. Therefore, 
$
  I_{N_2-j,j+1} \circ I_{N_1-j,j+1}^{-1} \circ I_{n_1-j,j+1} \circ I_{n_2-j,j+1}^{-1} (t) \geq t
$
whenever $\frac{n_1}{n_2}\leq 1$ and $\frac{N_2}{N_1} \leq 1$. Observe that
\begin{align*}
&I_{N_2-j,j+1} \circ I_{N_1-j,j+1}^{-1} \circ I_{n_1-j,j+1} \circ I_{n_2-j,j+1}^{-1} (t) \geq t, \\ 
&\quad\implies I_{n_1-j,j+1} \circ I_{n_2-j,j+1}^{-1} (t) \geq I_{N_1-j,j+1} \circ I_{N_2-j,j+1}^{-1}(t),   \\
&\quad\implies I_{n_1-j,j+1} \circ I_{n_2-j,j+1}^{-1} (x) - I_{N_1-j,j+1} \circ I_{N_2-j,j+1}^{-1}(t) \geq t.
\end{align*}
Technical condition \eqref{eq:technical-condition} holds and the conclusion of Theorem~\ref{theorem:samedistgeneral} applies.
\end{proof}



\vspace{12pt}

\textbf{We think that Proposition 3.5 can be generalized to $\frac{n_1}{N_1} \leq \frac{n_2}{N_2}$ with a different argument.}

\vspace{12pt}

We tried the following in an attempt to generalize: \\

\textcolor{red}{Let
$$
  f_{n_1,n_2,N_1,N_2,j}(x) = I_{n_1-j,j+1} \circ I_{n_2-j,j+1}^{-1}(x) - I_{N_1-j,j+1} \circ I_{N_2-j,j+1}^{-1}(x).
$$
Applications of the chain rule, the inverse function theorem \citep{rudin1964principles} on Page 221, and this property of regularized incomplete beta functions
$$
  I^\prime_{\alpha,\beta}(I^{-1}_{\alpha,\beta}(x)) = \frac{I^{-1}_{\alpha,\beta}(x)^{\alpha-1}(1 - I^{-1}_{\alpha,\beta}(x))^{\beta-1}}{B(\alpha,\beta)}, 
$$
allow us to write 
\begin{align*}
	&f^\prime_{n_1,n_2,N_1,N_2,j}(x) = \frac{ I^\prime_{n_1-j,j+1} \circ I_{n_2-j,j+1}^{-1}(x)}{ I^\prime_{n_2-j,j+1} \circ I_{n_2-j,j+1}^{-1}(x)} - \frac{ I^\prime_{N_1-j,j+1} \circ I_{N_2-j,j+1}^{-1}(x)}{ I^\prime_{N_2-j,j+1} \circ I_{N_2-j,j+1}^{-1}(x)} \\
	&\qquad= \frac{\left(I_{n_2-j, j+1}^{-1}(x)\right)^{n_1-j-1}\left(1 - I_{n_2-j, j+1}^{-1}(x)\right)^{j}/B(n_1-j,j+1)}{\left(I_{n_2-j, j+1}^{-1}(x)\right)^{n_2-j-1}\left(1 - I_{n_2-j, j+1}^{-1}(x)\right)^{j}/B(n_2-j,j+1)} \\
	&\qquad\qquad - \frac{\left(I_{N_2-j, j+1}^{-1}(x)\right)^{N_1-j-1}\left(1 - I_{N_2-j, j+1}^{-1}(x)\right)^{j}/B(N_1-j,j+1)}{\left(I_{N_2-j, j+1}^{-1}(x)\right)^{N_2-j-1}\left(1 - I_{N_2-j, j+1}^{-1}(x)\right)^{j}/B(N_2-j,j+1)}\\
	%&= \frac{B(n_2 - j, j+1)}{B(n_1-k,k+1)}g_{n_2, j}^{-1}(x)^{n_2 - n_1}(1-g_{n_2, j}^{-1}(x))^{k-j} - \frac{B(N_2 - j, j+1)}{B(N_1-k,k+1)}g_{N_2, j}^{-1}(x)^{N_2 - N_1}(1-g_{N_2, j}^{-1}(x))^{k-j}\\
	&\qquad= \frac{B(n_2-j,j+1)}{B(n_1-j,j+1)}\left(I_{n_2-j, j+1}^{-1}(x)\right)^{n_1 - n_2} - \frac{B(N_2 -j, j+1)}{B(N_1-j,j+1)}\left(I_{N_2-j, j+1}^{-1}(x)\right)^{N_1 - N_2} \\
	&\qquad= \left(I_{n_2-j,j+1}^{-1}(x)\right)^{n_1-n_2}\frac{B(N_2-j,j+1)}{B(N_1-j,j+1)} \\ 
	&\qquad\qquad \times\left[ \frac{B(n_2 -j, j+1)/B(n_1-j,j+1)}{B(N_2-j,j+1)/B(N_1-j,j+1)} - \frac{\left(I_{N_2-j, j+1}^{-1}(x)\right)^{N_1 - N_2}}{\left(I_{n_2-j, j+1}^{-1}(x)\right)^{n_1 - n_2}}\right]
\end{align*}
From here we see that $f_{n_1,n_2,N_1,N_2,j}(0) = f_{n_1,n_2,N_1,N_2,j}(1) = 0$, and when $x = 1$, we have 
\begin{align*}
	&f^\prime_{n_1,n_2,N_1,N_2,j}(1) = \frac{B(n_2-j,j+1)}{B(n_1-j,j+1)} - \frac{B(N_2 -j, j+1)}{B(N_1-j,j+1)}\\
	&\qquad= \frac{(n_1 - j)\dots(n_1)}{(n_2 - j)\dots(n_2)} - \frac{(N_1 - j)\dots(N_1)}{(N_2 - j)\dots(N_2)} \\
	&\qquad\leq 0,
\end{align*}
from the assumption that $\frac{n_1}{N_1} \leq \frac{n_2}{N_2}$. If we can prove the $\frac{\left(I_{N_2-j, j+1}^{-1}(x)\right)^{N_1 - N_2}}{\left(I_{n_2-j, j+1}^{-1}(x)\right)^{n_1 - n_2}}$ is a monotonic decreasing function for $x \in (0,1)$ and ranges from 0 to 1, then we can conclude that $f_{n_1,n_2,N_1,N_2,j}(x) \geq 0$ as desired.}



\bibliography{FHM.bib}
\bibliographystyle{plainnat}


\end{document}


